\chapter*{\chapterstyle{Analyse I {:} Propriétés de \(\R\)}}
On considère \(\R\) le corps des réels, et on le munit de sa \textbf{relation d'ordre totale} usuelle.\+
Soit \(A \subseteq \R\).

\subsection*{\subsecstyle{Majorants {:}}}
Cette relation d'ordre nous permet de parler des \textbf{majorants} d'une partie de \(\R\). Les propriétés sont identiques si on considère les \textbf{minorants} de \(\R\).

Soit \(m \in \R\), alors \(m\) est un \textbf{majorant} de \(A\) si il vérifie:
\[
    \boxed{
        \forall a \in A \; ; \; m \geq a
    }
\]
En outre, si \(m \in A\), on dit que \(m\) et le \textbf{maximum} de \(A\).\+
Soit \(M^+\) \textbf{l'ensemble des majorants} de la partie \(A\), alors on appelle \textbf{borne supérieure}, le nombre \(\sup(A)\) défini comme suit:
\[
    \boxed{
        \sup(A) := \min(M^+)
    }
\]
Soit \(m\) un majorant de \(A\) et \(\epsilon\) strictement positif, alors \(m\) est la borne supérieure de \(A\) si et seulement si il vérifie l'une des conditions équivalentes:

\[
    \begin{aligned}
        \framebox(140,15){
            \exists a \in A \; ; \; m - \epsilon \leq a \leq m
        }
    \end{aligned}
    \hspace{40pt}
    \begin{aligned}
        \framebox(140,15){
            \exists (u_n) \in A^\N \; ; \; u_n \longrightarrow m
        }
    \end{aligned}
\]
La propriété fondamentale de l'ensemble \(\R\) est que \textbf{toute partie non vide et majorée admet une borne supérieure}.

\subsection*{\subsecstyle{Propriété d'Archimède}}
Soit \(K\) un corps totalement ordonné, alors \(K\) est dit \textbf{archimédien}, si et seulement si on a:
\[
    \boxed{
        \forall (x, a) \in K_+ \; , \; \exists n \in \N \; ; \; na > x
    }
\]
Le cas trivial \(x < 0\) est exclu pour plus de lisibilité.\+
Informellement, \(K\) est archimédien, si tout nombre de \(K\) peut être ``dépassé`` par un multiple entier d'un autre nombre de \(K\).

\subsection*{\subsecstyle{Valeur absolue {:}}}
On introduit la fonction \textbf{valeur absolue} qui sera très utile par la suite pour mesurer des écarts, elle est définie comme suit:
\[
    |x| := \max(x, -x)
\]
Cette fonction vérifie plusieurs propriétés fondamentales, elle est \textbf{multiplicative} et vérifie les deux inégalités triangulaires:
\begin{flalign*}
    &\bullet |x + y| \leq |x| + |y| \shorteqnote{(Première inégalité triangulaire)} \\
    &\bullet ||x| - |y|| \leq |x - y| \shorteqnote{(Seconde inégalité triangulaire)}
\end{flalign*}
Ces inégalités se généralisent à des familles finies \(a_1, a_2, \ldots, a_n\), ie on a:
\[
    \Biggl|\sum_{i = 0}^n a_i\Biggl| \leq \sum_{i = 0}^n \bigl|a_i\bigl|
\]
\subsection*{\subsecstyle{Partie entière {:}}}
Soit \(n \in \N\), on appelle \textbf{partie entière} d'un réel \(x\), et on note \(\lfloor x \rfloor\) \textbf{l'unique} entier relatif tel que:
\[
    \lfloor x \rfloor \leq x < \lfloor x \rfloor + 1
\]
La partie entière vérifie la propriété fondamentale suivante:
\[
    \lfloor x + n \rfloor = \lfloor x \rfloor + n
\]
L'existence de la partie entière est garantie par \textbf{la propriété d'Archimède}.
\subsection*{\subsecstyle{Approximation décimale {:}}}
Soit \(n \in \N\), il peut être utile d'approximer un réel à \(n\) chiffres après la virgule, ie à \(10^{-n}\) près, la partie entière nous permet alors de le faire, et on l'obtient en exécutant l'algorithme suivant:
\[
    x \rightarrow 10^n x \rightarrow \lfloor 10^n x \rfloor \rightarrow \frac{\lfloor 10^n x \rfloor}{10^n}
\]
On obtient alors une approximation à \textbf{n chiffre après la virgule} de \(x\) qu'on note \(d_n(x)\) et on a:
\[
    \boxed{
        d_n(x) = \frac{\lfloor 10^n x \rfloor}{10^n}
    }
\]



\chapter*{\chapterstyle{Analyse I {:} Suites numériques}}
On appelle \textbf{suite} une \textbf{famille d'éléments} indexée par une partie \(I\) de \(\N\), et on note une telle famille \((u_n)_{n \in I}\), qu'on abrège souvent en \((u_n)\).\<

Lorsque tout les éléments de la famille appartiennent au \textbf{même ensemble} \(E\), on peut alors assimiler la suite \((u_n)\) à une application de \(\N\) vers \(E\).\+
Dans la suite on considérera uniquement des suites à valeurs dans un corps commutatif, noté \(K\).\<

Soit \((u_n)\), \((v_n)\) et \((w_n)\) trois suites à valeurs dans K.

\subsection*{\subsecstyle{Zoologie {:}}}
Il existe plusieurs types de suites aux propriétés singulières:
\begin{flalign*}
    &\bullet \; \text{Les suites \textbf{monotones}} \shorteqnote{(Ce sont les suite croissantes ou décroissantes)}\\
    &\bullet \; \text{Les suites \textbf{constantes}}\\ 
    &\bullet \; \text{Les suites \textbf{périodiques}}
\end{flalign*}
Un cas particulier étant les suites \textbf{stationnaires}, qui sont constantes \textbf{à partir d'un certain rang}.\<

Il faut noter que parler de suite \textbf{monotones} nécessite d'avoir une suite à valeurs dans un corps totalement ordonné.
Dans ce cas précis, on peut alors parler de suites \textbf{majorées ou minorées}. Dans le cas d'une suite à valeurs dans \(\C\), on peut néanmoins parler de suites bornées qui vérifient la proposition:
\[
    \exists M \in \R_+ \; , \; \forall n \in \N \; ; \; |u_n| \leq M
\]
Les suites à valeur dans \(\C\) vérifient la propriété fondamentale:
\[
    u_n = \Re(u_n) + i\Im(u_n)
\]

\subsection*{\subsecstyle{Convergence {:}}}
Si la suite \((u_n)\) vérifie la proposition ci-dessous, on dit qu'elle est \textbf{convergente} de limite \(l \in K\), sinon on dit qu'elle est \textbf{divergente}.
\[
    \boxed{
        \exists l \in K \; , \; \forall \epsilon \in \R_+ \; , \; \exists N \in \N \; , \; \forall n > N \; ; \; |u_n - l| < \epsilon   
    }
\]

Si la suite \((u_n)\) vérifie la proposition ci-dessous, on dit qu'elle tends vers \textbf{\(+\infty\)} (on définit de manière analogue les suites qui tendent vers \(-\infty\)):
\[
    \boxed{
        \forall M \in \R_+ \; , \; \exists N \in \N \; , \; \forall n > N \; ; \; u_n > M
    }
\]

Si une suite tends vers une limite \(l \in \overline{K}\), alors cette limite est \textbf{unique} et on note \(\)
\[
    u_n \rightarrow l
\]
On alors les propriétés fondamentales suivantes:
\begin{align*}
    &\bullet \text{Si \(u_n \rightarrow \pm\infty\), alors elle n'est pas majorée / minorée} \\
    &\bullet \text{Si \(u_n\) converge, alors elle est bornée}\\
    &\bullet \text{Si \(u_n \rightarrow l\), alors \(|u_n| \rightarrow |l|\)}
\end{align*}
Enfin, une suite à valeur dans \(\C\) converge si et seulement si \(\Re(u_n)\) et \(\Im(u_n)\) convergent.




\subsection*{\subsecstyle{Opérations sur les limites {:}}}
On considère que \((u_n)\) et \((v_n)\) sont convergentes de limites finies \(l, l' \in K^*\).\+
Soit \(\lambda \in K\), on a alors les propriétés suivantes:
\begin{align*}
    (u_n + v_n) &\rightarrow l + l' \\
    (u_n v_n) &\rightarrow ll' \\
    (\lambda u_n) &\rightarrow \lambda l
\end{align*}
Dans le cas où les limites sont infinies ou nulles, des \textbf{formes indéterminées} peuvent apparaître et nécessitent les précautions usuelles pour conclure.\<

Dans le cas d'un corps totalement ordonné, si il existe un \(m\) tel que pour tout \(n\), on ait \(u_n < m\) alors on peut en conclure:
\begin{flalign*}
        l \leq m
\end{flalign*}
Qui est analogue au \textbf{passage à la borne supérieure}.

\subsection*{\subsecstyle{Suites réelles {:}}}
On considère \((u_n)\), \((v_n)\) et \((w_n)\) des suites à valeurs dans \(\R\).
On a alors le \textbf{théorème de la limite monotone}:

\begin{center}
    \fbox{
        \parbox{0.771\linewidth}{%
            Si la suite \((u_n)\) est \textbf{croissante et majorée} alors elle converge vers \(\sup\{ u_n \; ; \; n \in \N\}\).\\
            Si la suite \((u_n)\) est \textbf{décroissante et minorée} alors elle converge vers \(\inf\{ u_n \; ; \; n \in \N\}\).\\
            Si la suite \((u_n)\) est \textbf{croissante} (resp décroissante) et \textbf{non majorée} (resp non minorée) alors elle tends vers respectivement \(\pm \infty\).
        }
    }
\end{center}

Supposons que \((u_n)\) et \((w_n)\) tendent vers une limite finie \(l \in K\), alors on a aussi le \textbf{théorème d'encadrement}:

\[
    \boxed{
            \Bigr[\forall n \in \N \; ; \; u_n \leq v_n \leq w_n \Bigr] \implies v_n \rightarrow l
    }
\]

On appelle \textbf{suites adjacentes} deux suites \((u_n)\) et \((v_n)\) telles que l'une est \textbf{croissante} et l'autre \textbf{décroissante} et telles que:
\[
    (u_n - v_n) \rightarrow 0
\]
On peut alors conclure que les deux suites convergent vers une limite \(l \in K\) et que cette limite est \textbf{encadrée} par ces deux suites.

\subsection*{\subsecstyle{Comparaison Asymptotique {:}}}
Si on veut comparer deux suites \((u_n)\) et \((v_n)\) pour des valeurs très grandes de \(n\), on parle alors de \textbf{comparaison asymptotique} et on définit plusieurs outils qui nous permettent de le faire:\<

On dit que la suite \((u_n)\) est \textbf{dominée} par la suite \((v_n)\) et on note alors \(u_n = O(v_n)\) si et seulement si il existe une suite \textbf{bornée} \((\theta_n)\) telle qu'à partir d'un certain rang, on ait la propriété ci dessous.\+
C'est une relation de \textbf{préordre} sur les suites, et elle est aussi \textbf{linéaire et multiplicative}.
\[
    u_n = \theta_n v_n
\]


On dit que la suite \((u_n)\) est \textbf{négligeable} devant la suite \((v_n)\) et on note \(u_n = o(v_n)\) si et seulement si il existe une suite \((\epsilon_n)\) qui tende vers 0, telle qu'à partir d'un certain rang, on ait la propriété ci dessous.\+
C'est une relation \textbf{transitive, linéaire et multiplicative}.
\[
    u_n = \epsilon_n v_n
\]

On dit que la suite \((u_n)\) est \textbf{équivalente} à la suite \((v_n)\) et on note \(u_n \sim v_n\) si et seulement si il existe une suite \((\eta_n)\) qui tende vers 1, telle qu'à partir d'un certain rang, on ait la propriété ci-dessous.\+ 
C'est une \textbf{relation d'équivalence} sur les suites, et elle est aussi \textbf{multiplicative}. 
\[
    u_n = \eta_n v_n
\]

\chapter*{\chapterstyle{Analyse I {:} Topologie de \(\R\)}}
Dans cette partie on admettra que \(\R\) est un \textbf{espace topologique}.

\subsection*{\subsecstyle{Adhérence {:}}}
On appelle \textbf{adhérence} d'une partie \(A\) de \(\R\) le plus petit fermé contenant \(A\). C'est un \textbf{opérateur de clôture} pour l'inclusion, ie c'est une application \textbf{idempotente, croissante, et extensive}.\<

Formellement, si on pose \((\mathcal{F}_i)_{i \in \N}\) la famille de tout les fermés contenant \(A\), on a:
\[
    \overline{A} := \bigcap_{i \in \N} \mathcal{F}_i
\]

\subsection*{\subsecstyle{Intérieur {:}}}
On appelle \textbf{intérieur} d'une partie \(A\) de \(\R\) le plus grand ouvert contenu dans \(A\). C'est une application \textbf{idempotente, croissante}.\<

Formellement, si on pose \((\mathcal{O}_i)_{i \in \N}\) la famille de tout les ouverts contenus dans \(A\), on a:
\[
    \accentset{\circ}{A} := \bigcup_{i \in \N} \mathcal{O}_i
\]
Il est intéressant de noter que l'adhérence et l'intérieur sont des notions duales.


\subsection*{\subsecstyle{Frontière {:}}}
On appelle \textbf{frontière} d'une partie \(A\) de \(\R\) et on note \(\partial A\) la partie qui vérifie une des conditions équivalentes ci-dessous:
\begin{align*}
    \partial A &= \overline{A} \backslash \accentset{\circ}{A} \\
    \partial A &= \overline{A} \cap \overline{E \backslash A} \\    
\end{align*}
Intuitivement, la frontière est \textbf{l'ensemble des points situés au ``bord`` de cet ensemble}.\+
En particulier, un ensemble est un ouvert si et seulement s'il est disjoint de sa frontière.

\subsection*{\subsecstyle{Densité {:}}}
La densité est une propriété fondamentale d'une partie d'une espace \textbf{topologique}.
Soit \(A \subseteq \R\), alors on dit que \(A\) est \textbf{dense dans} \(\R\) si et seulement si elle vérifie une des conditions équivalentes ci-dessous:
\begin{center}
    \fbox{%
        \parbox{0.42\linewidth}{%
            &\bullet \; \text{Tout \textbf{ouvert} de \(\R\) contient des éléments de \(A\)}.\\
            &\bullet \, \text{L'adhérence de \(A\) est égale à \(\R\).}
        }
        %
    }
\end{center}

En particulier, on a par exemple les exemples canoniques suivants:
\begin{align*}
    &\bullet \text{\(\R\) est dense dans lui-même}.\\
    &\bullet \text{\(\Q\) est dense dans \(\R\)}.\\
    &\bullet \text{\(\R \backslash \Q\) est dense dans \(\R\)}.
\end{align*}

\subsection*{\subsecstyle{Extractions {:}}}
On appelle extraction toute fonction \(\varphi : \N \rightarrow \N\) \textbf{strictement croissante}.\+
On dit que \(v_n\) est une \textbf{suite extraite} de \(u_n\) si il existe une extraction \(\varphi\) telle que:
\[
    v_n = u_{\varphi(n)}
\]
On note aussi qu'une composée d'extraction est encore une extraction.
\subsection*{\subsecstyle{Valeurs d'adhérence {:}}}
On appelle \textbf{valeur d'adhérence} tout limite d'une suite-extraite de \(u_n\) et on peut caractériser ces valeurs d'adhérence comme suit:
\[
    \forall \epsilon \in \R_+ \; , \; \forall N \in \N \; , \; \exists n > N \; ; \; |u_n - l| < \epsilon 
\]
On peut alors montrer cette propriété fondamentale:
\[
    \Bigr[u_n \rightarrow l \in \overline{K} \Bigr] \implies \Bigr[ \forall \varphi \; ; \; u_{\varphi(n)} \rightarrow l \Bigr]
\]
On montre aussi le \textbf{théorème de Bolzano-Weirstrass}:
\begin{center}
    \fbox{
        \textbf{Tout suite bornée admet au moins une valeur d'adhérence.}
    }
\end{center}

\subsection*{\subsecstyle{Complétude {:}}}
On appelle \textbf{suite de Cauchy} toute suite qui vérifie:
\[
    \forall \epsilon \in \R_+ \; , \; \exists N \in \N \; , \; \forall p, q > N \; ; \; |u_p - u_q| < \epsilon 
\]
Informellement, on en déduit que les termes de la suite sont très proches \textbf{les uns des autres} à partir d'un certain rang.\+
\textbf{Toute suite convergente} est de Cauchy.\<

Supposons \(u_n\) une suite de Cauchy, alors elle vérifie les propriétés suivantes:
\begin{align*}
    &\bullet u_n \text{ est bornée.} \\
    &\bullet u_n \text{ converge si elle admet une valeur d'adhérence.}
\end{align*}

On appelle \textbf{espace métrique} tout espace topologique muni d'une \textbf{distance}, par exemple \((\R, |\cdot|)\), ou \((\C, |\cdot|)\).\+
On appelle \textbf{espace complet} tout espace métrique dans lequel \textbf{tout suite de Cauchy} converge.\<

En particulier \(\R\) et \(\C\) sont des espaces complets, et tout sous-ensemble fermé d'un espace complet et aussi complet. \<

